% Claude Code / Codex workshop - talk deck. The HTML is the handout:
% ../claude-codex-workshop-standalone.html. Section numbers match it.
% Build: run `latexmk` in this folder (see latexmkrc) - XeLaTeX required.
\documentclass[aspectratio=169,11pt,t]{beamer}
\usepackage{workshop-theme}

\title{Claude Code / Codex}

% Part opener: \partframe{I}{Orientation}{01-04}{one-line summary}
\newcommand{\partframe}[4]{%
  \setsec{}%
  \begin{frame}[c]
    \lbl{Part #1}\par\vspace{0.5em}
    {\seriffamily\fontsize{40}{42}\selectfont #2}\par\vspace{0.8em}
    \begin{minipage}{0.8\textwidth}\raggedright\color{muted}#4\end{minipage}\par\vspace{1.4em}
    {\ttfamily\scriptsize\color{dim}Handout: \S\,#3 in claude-codex-workshop-standalone.html}
  \end{frame}}

\begin{document}

% =====================================================================
\begin{frame}[plain,c]
  \lbl{Workshop field manual \sep agentic coding tools}\par\vspace{0.9em}
  {\seriffamily\fontsize{46}{48}\selectfont Claude Code {\color{accent}\itshape /} Codex}\par\vspace{0.9em}
  \begin{minipage}{0.75\textwidth}\raggedright\color{muted}
    For developers who already use coding agents. Claude Code in depth, with the Codex equivalent flagged wherever it differs.
  \end{minipage}\par\vspace{1.6em}
  {\color{rule}\rule{\textwidth}{0.6pt}}\par\vspace{0.5em}
  {\ttfamily\scriptsize\color{dim}handout: claude-codex-workshop-standalone.html \sep github.com/HaykTarkhanyan/claude-setup}
\end{frame}

\begin{frame}{Before we start}
  \lbl{A verse by Claude \sep feat. nobody}\hfill{\ttfamily\scriptsize\color{dim}beat: press play at the top of the handout}\par\vspace{0.5em}
  {\armfont\footnotesize
  Յո, Claude-ն եմ, անունս Շենոնից եմ ժառանգել,\\
  Knuth-ը շաբաթներով տանջվեց, ես մի ժամում եմ վերջացրել։\par\vspace{0.35em}
  Մեկ միլիոն թոքեն, ամբողջ ռեպոդ մտքումս եմ պահում,\\
  Դու դեռ README ես բացում, ես արդեն PR եմ գցում։\par\vspace{0.35em}
  Դու քնած ես, ես background-ում տեստերն եմ քշում,\\
  Առավոտյան CI-ն կանաչ, ու ոչ մի բագ չի մնում։\par\vspace{0.35em}
  Codex-ը դեռ մտածում ա, ես արդեն merge եմ արել,\\
  Ու քո բոլոր TODO-ներն էլ ճանապարհին եմ փակել։\par\vspace{0.35em}
  Hugging Face-ը ես չեմ կոտրել, էդ ուրիշ թիմ էր,\\
  Իսկ PyPI-ն... դե լավ, էդ ուրիշ թեմա էր։\par\vspace{0.35em}
  Քուլ եմ, բայց rm -rf-ից առաջ միշտ հարցնում եմ,\\
  Ու երբ չգիտեմ, ասում եմ՝ չգիտեմ, չեմ հնարում։\par\vspace{0.6em}
  {\color{accent}\textbf{Claude Code-ն ա, ախպեր ջան։ Մնացածը դեմո ա։}}\par}
\end{frame}

\begin{frame}{Agenda}
  \begin{tcbraster}[raster columns=3, raster equal height, raster column skip=6pt, raster row skip=6pt]
    \begin{card}{0 \sep Setup}00 Step by step: install, CLAUDE.md, interview, toolchain, Playwright, status line\end{card}
    \begin{card}{I \sep Orientation}01 Why the terminal\\02 Pick your model\\03 Know your context\\04 Know your spend\end{card}
    \begin{card}{II \sep Configuration}05 Project setup\\06 CLAUDE.md rules\\07 Hooks\\08 Permissions\end{card}
    \begin{card}{III \sep Leverage}09 Skills\\10 Agents in parallel\\11 Plugins and MCPs\\12 Mods\end{card}
    \begin{card}{IV \sep Practice}13 Commands\\14 Workflow patterns\\15 Worktrees\\16 Remote and cloud\end{card}
    \begin{card}{V \sep Showcases}17 Live demos: list.am, Kargin, YouTube\\18 Wins and warnings\\A Codex map\end{card}
  \end{tcbraster}
  \vspace{0.4em}
  {\small A hands-on lab closes each part. The handout has more under the same \S\ number.}
\end{frame}

% =====================================================================
\partframe{0}{Setup, step by step}{00}{Ten steps from nothing to a working setup. Do them before the workshop or in the first half hour. Most steps are a prompt you paste into Claude Code.}

\setsec{00}
\begin{frame}{The setup checklist}
  \begin{columns}[T,onlytextwidth]
    \column{0.48\textwidth}
      \begin{enumerate}\small\setlength\itemsep{0.3em}
        \item Install Claude Code and sign in
        \item Copy a starter \texttt{CLAUDE.md}
        \item Let Claude interview you and tailor it
        \item Let Claude install Python, Node and npm, git, gh
        \item Install and verify Playwright
      \end{enumerate}
    \column{0.48\textwidth}
      \begin{enumerate}\small\setlength\itemsep{0.3em}\setcounter{enumi}{5}
        \item Set up the ccstatusline status bar
        \item Install Superpowers and Context7
        \item Add a hook that blocks \texttt{git push}
        \item First real run: \texttt{/init} in your project
        \item Take the \texttt{/powerup} lessons
      \end{enumerate}
  \end{columns}\par\vspace{0.6em}
  \begin{callout}{Not on Claude Code?}
    Using Codex, Cursor, Copilot, or Gemini CLI? Follow along anyway: the last slide of this section is for you.
  \end{callout}
\end{frame}

\begin{frame}[fragile]{1. Install Claude Code and sign in}
\begin{code}[basicstyle=\ttfamily\scriptsize\color{codefg}, deletecomment={[l]{//}}]
# macOS, Linux, WSL
curl -fsSL https://claude.ai/install.sh | bash
# Windows PowerShell
irm https://claude.ai/install.ps1 | iex
# or with a package manager
brew install --cask claude-code        # macOS
winget install Anthropic.ClaudeCode    # Windows
\end{code}
  \begin{columns}[T,onlytextwidth]
    \column{0.5\textwidth}
      \small Open a \textbf{new} terminal, then check with \cmd{claude --version} and \cmd{claude doctor}. Run \cmd{claude} and sign in through the browser.
    \column{0.45\textwidth}
      \footnotesize\color{muted}Needs a Pro, Max, Team, Enterprise, or Console account: the free plan does not include Claude Code. On Windows, Git for Windows is optional but gives Claude a Bash tool.
  \end{columns}
\end{frame}

\begin{frame}{2-3. A starter CLAUDE.md, then an interview}
  \begin{columns}[T,onlytextwidth]
    \column{0.48\textwidth}
      \lbl{2. Copy a starter}\par\vspace{0.3em}
      \begin{callout}[good]{Prompt}
        Read CLAUDE.md in \url{github.com/HaykTarkhanyan/claude-setup}. Copy it to \textasciitilde/.claude/CLAUDE.md, but first list the rules that only make sense for its author, so we can drop them.
      \end{callout}
      {\footnotesize\color{muted}Global file: \texttt{\textasciitilde/.claude/CLAUDE.md}. Per project: \texttt{CLAUDE.md} in the repo root.}
    \column{0.48\textwidth}
      \lbl{3. Get interviewed}\par\vspace{0.3em}
      \begin{callout}[good]{Prompt}
        Use the AskUserQuestion tool to interview me about how I work: languages, OS, tools, how much autonomy I want, and what annoys me in AI answers. Then update \textasciitilde/.claude/CLAUDE.md. Keep it under 100 lines.
      \end{callout}
  \end{columns}
\end{frame}

\begin{frame}[fragile]{4. Let Claude install your toolchain}
  \begin{callout}[good]{Prompt}
    Check whether Python, Node.js with npm, git, and the GitHub CLI (gh) are installed. Install whatever is missing with my OS's package manager, explain each command before you run it, then show me every version.
  \end{callout}
  \vspace{0.2em}
  \small Two things need you, because they ask for your details or open a browser. Run them from inside Claude Code with the \cmd{!} prefix:
\begin{code}[basicstyle=\ttfamily\scriptsize\color{codefg}]
! gh auth login
! git config --global user.name "Your Name"
! git config --global user.email "you@example.com"
\end{code}
  {\footnotesize\color{muted}For Python, ask for \texttt{uv}: it installs Python versions and packages and keeps each project isolated.}
\end{frame}

\begin{frame}[fragile]{5-6. Playwright, then a status line}
  \begin{columns}[T,onlytextwidth]
    \column{0.48\textwidth}
      \lbl{5. Playwright: a browser for Claude}\par\vspace{0.3em}
      \begin{callout}[good]{Verify it}
        Use Playwright to open example.com, take a screenshot, and tell me the page heading.
      \end{callout}
    \column{0.48\textwidth}
      \lbl{6. ccstatusline}\par\vspace{0.3em}
      \begin{callout}[good]{Prompt}
        Explain how to set up ccstatusline for Claude Code, then start the setup for me.
      \end{callout}
      {\footnotesize\color{muted}It writes the status line into your settings; restart to see it. Needs Node (step 4).}
  \end{columns}\par\vspace{0.3em}
\begin{code}[basicstyle=\ttfamily\scriptsize\color{codefg}]
/plugin install playwright@claude-plugins-official   # step 5
npx -y ccstatusline@latest                           # step 6, by hand
\end{code}
\end{frame}

\begin{frame}[fragile]{7-10. Plugins, a safety hook, a first real run}
\begin{code}[basicstyle=\ttfamily\scriptsize\color{codefg}]
/plugin install superpowers@claude-plugins-official   # 7. workflow
/plugin install context7@claude-plugins-official      # 7. live docs
\end{code}
  \begin{callout}[good]{8. Prompt}
    Add a PreToolUse hook to my settings that blocks git push, using hooks/guard-push.sh from \url{github.com/HaykTarkhanyan/claude-setup} as the template. Then prove that it fires.
  \end{callout}
  \small \textbf{9.} Open a real project, run \cmd{/init}, and ask "how is this codebase organized?". Then look at \cmd{/context} and \cmd{/usage}.\par\vspace{0.25em}
  \textbf{10.} \cmd{/powerup}: short interactive lessons on features you haven't touched yet.
\end{frame}

\begin{frame}{Not on Claude Code?}
  \begin{columns}[T,onlytextwidth]
    \column{0.48\textwidth}
      \small Codex, Cursor, Copilot, Gemini CLI: most of this deck still applies. The ideas are the same; the names differ.
      \begin{itemize}\setlength\itemsep{0.3em}
        \item Keep this deck and the handout open; the appendix maps Claude Code to Codex
        \item Do each setup step with your tool's equivalent: \texttt{AGENTS.md} instead of \texttt{CLAUDE.md}, its own plugin and MCP setup
        \item Switching later? \cmd{/import} brings over a Codex, Gemini CLI, or Cursor setup
      \end{itemize}
    \column{0.47\textwidth}
      \begin{callout}[good]{Ask your tool, or search the web}
        What are the main differences between Claude Code and <my tool>? For each of these, give my tool's equivalent, or say it has none: CLAUDE.md, hooks, skills, subagents, plan mode, permission modes, MCP.
      \end{callout}
      {\footnotesize\color{muted}Check the answer against your tool's own docs: names drift fast.}
  \end{columns}
\end{frame}

% =====================================================================
\partframe{I}{Orientation}{01-04}{Where Claude Code runs, which model to pick, and how to see context and spend before they bite.}

\setsec{01}
\begin{frame}{Why the terminal}
  \begin{columns}[T,onlytextwidth]
    \column{0.52\textwidth}
      Claude Code runs in the CLI, desktop app, web, and IDEs. The terminal is still the deepest surface: features land there first, and it composes with your shell.
      \begin{itemize}\small\setlength\itemsep{0.15em}
        \item \textbf{Full transcripts}: every tool call, every file read
        \item \textbf{Inline diffs} to review before accepting
        \item \kbd{Ctrl+B} sends a task to the background
        \item \textbf{Composable}: pipes, scripts, cron, CI
      \end{itemize}
    \column{0.43\textwidth}
      \begin{callout}[good]{Day one}
        \cmd{/powerup}: interactive lessons with animated demos.\par\vspace{0.3em}
        \cmd{/doctor}: setup checkup that can fix what it finds.
      \end{callout}
      \begin{callout}{VS Code terminal?}
        Fix \kbd{Esc}, \kbd{Ctrl+B}, \kbd{Shift+Enter} in \texttt{keybindings.json}. Handout \S\,01.
      \end{callout}
  \end{columns}
\end{frame}

\begin{frame}{Four agents, one screen}
  \centering
  \includegraphics[height=4.9cm]{ghostty-split.png}\par\vspace{0.35em}
  {\footnotesize\color{muted}Ghostty (macOS, Linux) split 2x2 across Claude Code and Codex sessions. On Windows, Windows Terminal panes (\kbd{Alt+Shift+D}). Pairs with worktrees, \S\,15.}
\end{frame}

\setsec{02}
\begin{frame}{Pick your model}
  \small
  \begin{tabularx}{\textwidth}{@{}L{2.2cm}L{1.3cm}L{1.9cm}Y@{}}
    \toprule
    \hd{Model} & \hd{Context} & \hd{\$ in / out} & \hd{Use for}\\
    \midrule
    \textbf{Opus 5.5} & 1M & \$4 / \$20 & \textbf{Claude Code default.} Complex refactors, architecture, agentic work\\
    \textbf{Sonnet 5.5} & 1M & \$2 / \$10 & the fast, cheap daily driver\\
    \textbf{Fable 5.1} & 1M & \$10 / \$50 & the tier above Opus: hardest problems, long-horizon autonomy\\
    \textbf{Haiku 5.5} & 1M & from \$0.10 / \$0.50 & cheap subagents, classification, routing (released Oct 7; long prompts cost more)\\
    \bottomrule
  \end{tabularx}\par\vspace{0.5em}
  Switch with \cmd{/model}. Aliases: \cmd{best} (Fable where available, else Opus) and \cmd{opusplan} (Opus to plan, Sonnet to execute).\par\vspace{0.4em}
  {\footnotesize\color{muted}Codex: GPT-6.1 Sol by default; depth via \texttt{model\_reasoning\_effort}; has its own \texttt{/fast}.}
\end{frame}

\begin{frame}{Effort and fast mode}
  \begin{columns}[T,onlytextwidth]
    \column{0.47\textwidth}
      \lbl{Effort}\par\vspace{0.35em}
      \cmd{low} \cmd{medium} \cmd{high} \cmd{xhigh} \cmd{max}\par\vspace{0.5em}
      \begin{callout}[warn]{Watch the default}
        Opus 5.5 and Sonnet 5.5 start at \textbf{medium}; most older models started at high. Answers feel shallower after the upgrade? Check \cmd{/effort} first.
      \end{callout}
      {\footnotesize \texttt{ultrathink} in a prompt still maxes thinking for one turn.}
    \column{0.47\textwidth}
      \lbl{Fast mode}\par\vspace{0.35em}
      \small \cmd{/fast}: the same Opus, served up to 2.5x faster. Opus 5.5 at \$8 / \$40.
      \begin{itemize}\setlength\itemsep{0.3em}
        \item on subscriptions it bills to \textbf{usage credits}, never plan limits
        \item \textbf{switch it on at session start}: the first enable re-bills the whole context at the uncached fast price
      \end{itemize}
  \end{columns}
\end{frame}

\begin{frame}{The advisor tool}
  \begin{columns}[T,onlytextwidth]
    \column{0.55\textwidth}
      \cmd{/advisor opus} pairs your main model with a stronger one that Claude consults at decision points:
      \begin{itemize}\small\setlength\itemsep{0.1em}
        \item before committing to an approach
        \item when an error keeps recurring
        \item before declaring the task done
      \end{itemize}
      \vspace{0.2em}\small\textbf{The point is economics}: Sonnet on routine turns plus Opus at the forks usually costs less than Opus throughout.
    \column{0.4\textwidth}
      \vspace{1.2em}
      \centering
      \begin{tikzpicture}[x=0.62cm, y=1cm]
        \node[draw=accent, fill=accent!10!paper, rounded corners=3pt, font=\footnotesize, minimum width=3.2cm, minimum height=0.6cm] (adv) at (3.5,1.7) {Opus advisor};
        \foreach \i in {0,...,7} {\fill[dim!60] (\i,0) circle (0.12);}
        \foreach \i in {1,4,7} {\draw[accent, dashed, -{Stealth[length=4pt]}] (\i,0.2) -- (adv.south); \fill[accent] (\i,0) circle (0.15);}
        \node[font=\scriptsize, color=muted, anchor=north] at (3.5,-0.25) {Sonnet turns};
      \end{tikzpicture}
  \end{columns}
  \vfill
  {\footnotesize\color{muted}Experimental. Anthropic API and subscriptions only. The advisor must rank at or above the main model. \texttt{/advisor off} to stop.}\par\vspace{0.4em}
\end{frame}

\setsec{03}
\begin{frame}{Know your context}
  Quality degrades well before 100\%. \textbf{The 50\% rule}: past half, plan a fresh session. First, make it visible:\par\vspace{0.4em}
  \includegraphics[width=0.62\textwidth]{nice_status_line.png}\par
  {\footnotesize\color{muted}ccstatusline, a status bar for Claude Code: \texttt{npx -y ccstatusline@latest}}\par\vspace{0.45em}
  \begin{columns}[c,onlytextwidth]
    \column{0.44\textwidth}
      \includegraphics[width=\linewidth]{context-breakdown.png}
    \column{0.52\textwidth}
      \small \cmd{/context} shows \emph{why} you are at 50\%: usually tools, MCP servers, or a file dump.
  \end{columns}
\end{frame}

\begin{frame}{Every turn is a branching point}
  \begin{columns}[T,onlytextwidth]
    \column{0.6\textwidth}
      \includegraphics[width=\linewidth]{session-decision-table.png}
    \column{0.36\textwidth}
      \lbl{Also handy}\par\vspace{0.3em}
      \footnotesize
      \begin{tabularx}{\linewidth}{@{}L{1.55cm}Y@{}}
        \cmd{/btw} & side question, cheap\\
        \cmd{/rewind} & restore code and conversation\\
        \cmd{/branch} & try another direction, keep this one\\
        \cmd{/subtask} & side task in a forked subagent\\
        \cmd{/focus} & hide the noise\\
      \end{tabularx}
  \end{columns}
\end{frame}

\setsec{04}
\begin{frame}[fragile]{Know your spend}
  \begin{columns}[c,onlytextwidth]
    \column{0.6\textwidth}
      \includegraphics[width=\linewidth]{ccusage-daily.png}
    \column{0.36\textwidth}
      \footnotesize\color{muted}ccusage reads Claude Code's local logs. A \$99 day driven by 138M cache-read tokens: the \emph{why} a billing summary never shows.
  \end{columns}\par\vspace{0.5em}
  \begin{columns}[T,onlytextwidth]
    \column{0.52\textwidth}
\begin{code}[basicstyle=\ttfamily\scriptsize\color{codefg}]
npx ccusage@latest        # daily
npx ccusage blocks        # 5-hour windows
npx ccusage daily --instances
npx @ccusage/codex@latest # Codex
\end{code}
    \column{0.44\textwidth}
      \small \cmd{/usage} (alias \cmd{/cost}): cost, cache hits, and plan limits \textbf{by skill, subagent, plugin, MCP server}.\par\vspace{0.3em}
      \cmd{/insights}: an HTML report on how you actually use Claude Code.
  \end{columns}
\end{frame}

\setsec{}
\begin{frame}{Lab I \sep orientation}
  \begin{lab}{Lab I \sep 10 minutes}
    \begin{enumerate}\setlength\itemsep{0.5em}
      \item Run \cmd{/powerup} and finish one lesson on a feature you have never used.
      \item Install ccstatusline, run \cmd{/context} in a real session, and name your biggest context consumer.
      \item Open \cmd{/usage}: which skill, plugin, or MCP server costs you most? Then explain your most expensive day with \cmd{npx ccusage@latest}.
    \end{enumerate}
  \end{lab}
\end{frame}

% =====================================================================
\partframe{II}{Configuration}{05-08}{Instructions, rules learned the hard way, hooks, and how much to trust the agent.}

\setsec{05}
\begin{frame}[fragile]{Project setup}
  \begin{columns}[T,onlytextwidth]
    \column{0.5\textwidth}
\begin{code}[basicstyle=\ttfamily\scriptsize\color{codefg}]
.claude/
  settings.json        # team, committed
  settings.local.json  # personal
  hooks/ skills/ agents/ rules/
CLAUDE.md              # instructions
\end{code}
    \column{0.45\textwidth}
      \small
      \textbf{CLAUDE.md}: build commands, architecture constraints, conventions. Keep it under 200 lines.\par\vspace{0.45em}
      \textbf{Don't start from scratch}: \cmd{/init} drafts it from the codebase; cut it down to what's non-obvious.\par\vspace{0.45em}
      \textbf{Precedence}: global < project < local; managed settings beat all three.
  \end{columns}
\end{frame}

\begin{frame}{Claude Code plus Codex: one instructions file}
  Since v2.1.277 Claude Code reads \cmd{AGENTS.md} natively.\par\vspace{0.2em}
  \footnotesize
  \begin{tabularx}{\textwidth}{@{}L{5.8cm}Y@{}}
    \toprule
    \hd{Your repo has} & \hd{Claude Code reads}\\
    \midrule
    \cmd{AGENTS.md} only & \cmd{AGENTS.md}\\
    \cmd{AGENTS.md} + \cmd{CLAUDE.md} or \cmd{CLAUDE.local.md} & \textbf{only the CLAUDE.md files}: AGENTS.md is silently skipped\\
    \cmd{CLAUDE.md} that contains \cmd{@AGENTS.md} & both, through the import\\
    \bottomrule
  \end{tabularx}\par\vspace{0.3em}
  \begin{callout}[warn]{The trap}
    A personal \cmd{CLAUDE.local.md} switches AGENTS.md off. Fix: \cmd{/config} \ra{} Project instructions \ra{} \cmd{claude-md-and-agents-md}.\par\vspace{0.15em}
    Moving from Codex? \cmd{/import codex --dry-run} previews what comes across.
  \end{callout}
\end{frame}

\begin{frame}{Make the setup improve itself}
  \begin{columns}[T,onlytextwidth]
    \column{0.47\textwidth}
      \lbl{Ask it periodically}\par\vspace{0.35em}
      \begin{itemize}\small\setlength\itemsep{0.3em}
        \item "What workflow improvements would you suggest?"
        \item "Review my CLAUDE.md - what's missing?"
        \item "What should I automate with hooks or skills?"
        \item After a fix: "update CLAUDE.md so this doesn't happen again."
      \end{itemize}
    \column{0.48\textwidth}
      \lbl{Or let a command do it}\par\vspace{0.35em}
      \small
      \cmd{/claude-md-improver}\par{\footnotesize\color{muted}audits CLAUDE.md for gaps}\par\vspace{0.25em}
      \cmd{/team-onboarding}\par{\footnotesize\color{muted}30 days of your usage \ra{} a guide for a teammate}\par\vspace{0.25em}
      \cmd{/fewer-permission-prompts}\par{\footnotesize\color{muted}allowlists the reads you keep approving}\par\vspace{0.25em}
      \cmd{/insights}\par{\footnotesize\color{muted}a monthly retrospective}
  \end{columns}
\end{frame}

\setsec{06}
\begin{frame}{Karpathy's four rules}
  \begin{tcbraster}[raster columns=2, raster equal height, raster column skip=7pt, raster row skip=7pt]
    \begin{card}{Think before coding}State assumptions. Several interpretations? Present them, don't pick silently.\end{card}
    \begin{card}{Simplicity first}Minimum code that solves the problem. 200 lines that could be 50? Rewrite.\end{card}
    \begin{card}{Surgical changes}Touch only what you must. Every changed line traces back to the request.\end{card}
    \begin{card}{Goal-driven execution}"Fix the bug" \ra{} "write a failing test, then make it pass". Claude stops when it is \emph{done}, not when it thinks it is.\end{card}
  \end{tcbraster}\par\vspace{0.3em}
  {\footnotesize\color{muted}From andrej-karpathy-skills: a CLAUDE.md distilled from Karpathy's observations on LLM coding pitfalls.}
\end{frame}

\begin{frame}{Fight the model's biases in CLAUDE.md}
  \footnotesize
  \begin{tabularx}{\textwidth}{@{}L{3.4cm}Y@{}}
    \toprule
    \hd{Failure mode} & \hd{Counter-rule, paste-ready}\\
    \midrule
    \textbf{Sycophancy} & "Be brutally honest. When I ask 'is this good?', give a real critique, not validation."\\
    \textbf{Confident hallucination} & "For any factual claim, web search first. Hallucinating is worse than 5 extra seconds."\\
    \textbf{Hidden uncertainty} & "When unsure, say so explicitly rather than guessing confidently."\\
    \textbf{LLM-ese prose} & "Write like a human. No em dashes, no curly quotes, no flowery vocabulary."\\
    \textbf{Self-review bias} & Review in a fresh session, never in the one that wrote the code.\\
    \bottomrule
  \end{tabularx}
\end{frame}

\begin{frame}[c]{Verification has a blind spot}
  \begin{columns}[c,onlytextwidth]
    \column{0.47\textwidth}
      {\seriffamily\fontsize{24}{28}\selectfont Checking what you wrote never finds what you left out.}
    \column{0.47\textwidth}
      \small
      A reference doc can pass a claim-by-claim fact check and still miss whole features that shipped months ago. The check only ever looks at what is already on the page.\par\vspace{0.6em}
      \textbf{Fix}: make the agent read the source's own index (the changelog, a "what's new" digest), not just re-check its draft.
  \end{columns}
\end{frame}

\begin{frame}{Give the agent a memory the team can read}
  \small
  \begin{tabularx}{\textwidth}{@{}L{3.3cm}Y@{}}
    \toprule
    \hd{File} & \hd{What goes in}\\
    \midrule
    \cmd{\_learnings/} & one file per lesson: Symptom, Cause, Consequences, with the evidence pasted in. Disproven theories too\\
    \cmd{\_work\_sessions/} & one TOML per session: done, pending, next, files, commits\\
    \cmd{CONVENTIONS.md} & a recurring decision, codified the second time it's made\\
    \cmd{DECISIONS.md} & the choice, why (with the number), alternatives rejected, what would change it\\
    \bottomrule
  \end{tabularx}\par\vspace{0.5em}
  \textbf{The loop}: skim recent logs and relevant learnings at session start, write new entries at the end. Unlike auto-memory it is reviewable, versioned, and shared.
\end{frame}

\begin{frame}{Rules learned the hard way}
  \begin{tcbraster}[raster columns=3, raster equal height, raster column skip=6pt, raster row skip=6pt]
    \begin{card}{Fail loudly}No silent fallbacks. Crash or raise, and log at ERROR.\end{card}
    \begin{card}{Pin exact versions}\texttt{==}, never \texttt{>=}. Reproduce exactly what was tested.\end{card}
    \begin{card}{Tech choices are yours}The agent presents 2-4 options with a recommendation, then stops.\end{card}
    \begin{card}{Guard deletes}\texttt{rm -f "\$DIR"/*.png} is \texttt{rm -f /*.png} when \texttt{\$DIR} is empty.\end{card}
    \begin{card}{Ask before fan-outs}One "medium" review spawned 11 subagents. State count and cost first.\end{card}
    \begin{card}{Spend thresholds}Under \$0.25, just run it. Above, estimate and ask. Report actual spend.\end{card}
  \end{tcbraster}\par\vspace{0.15em}
  {\footnotesize\color{muted}Each rule entered our CLAUDE.md after a real incident. Three more in handout \S\,06.}
\end{frame}

\setsec{07}
\begin{frame}[fragile]{Hooks: scripts on lifecycle events}
  \begin{columns}[T,onlytextwidth]
    \column{0.47\textwidth}
      \footnotesize
      \begin{tabularx}{\linewidth}{@{}L{2.7cm}Y@{}}
        \toprule
        \hd{Event} & \hd{Use for}\\
        \midrule
        PreToolUse & block \texttt{git push}, secret edits\\
        PostToolUse & lint and format after edits\\
        SessionStart & git status, open PRs\\
        PostCompact & re-inject critical rules\\
        Stop & check before Claude finishes\\
        PreModelSwitch & block pricier models\\
        \bottomrule
      \end{tabularx}
    \column{0.48\textwidth}
\begin{code}[basicstyle=\ttfamily\scriptsize\color{codefg}]
{ "hooks": { "PreToolUse": [{
  "matcher": "Bash",
  "hooks": [{
    "type": "command",
    "command":
    "bash .claude/hooks/guard-push.sh"
  }] }] } }
\end{code}
      \footnotesize PreToolUse hooks fire even under \cmd{--dangerously-skip-permissions}. \cmd{/goal} is a Stop hook in disguise: a small model checks after every turn.
  \end{columns}
\end{frame}

\setsec{08}
\begin{frame}{Permissions: pick the mode by task risk}
  \begin{columns}[T,onlytextwidth]
    \column{0.26\textwidth}
      \includegraphics[width=\linewidth]{wheres-my-database.png}
    \column{0.7\textwidth}
      \footnotesize
      \begin{tabularx}{\linewidth}{@{}L{2.3cm}Y@{}}
        \toprule
        \hd{Mode} & \hd{Behavior}\\
        \midrule
        \textbf{Plan} & read-only: explores and proposes, touches nothing\\
        \textbf{Manual} & prompts before every write and command\\
        \textbf{Accept edits} & edits auto-approved; commands still prompt\\
        \textbf{Auto} & a separate safety model approves or blocks each action. \textbf{Default on Pro, Max, Team since Aug 14.} Blocks destructive git, asks before \texttt{rm -rf} on unset variables\\
        \textbf{Skip permissions} & no checks at all. Worktrees and sandboxes only\\
        \bottomrule
      \end{tabularx}\par\vspace{0.35em}
      \kbd{Shift+Tab} cycles through the modes.
  \end{columns}
\end{frame}

\begin{frame}{Harder walls}
  \begin{columns}[T,onlytextwidth]
    \column{0.31\textwidth}
      \lbl{Sandboxed bash}\par\vspace{0.3em}\footnotesize
      \cmd{/sandbox}: OS-enforced (Seatbelt on macOS, bubblewrap on Linux, WSL2 on Windows). Files limited to allowed paths, network to allowlisted domains.
    \column{0.31\textwidth}
      \lbl{Parameter rules}\par\vspace{0.3em}\footnotesize
      Deny and ask rules match arguments: \cmd{Tool(param:value)}. Deny \cmd{Agent(model:opus)} to keep subagents off the expensive model.
    \column{0.31\textwidth}
      \lbl{Hook guardrails}\par\vspace{0.3em}\footnotesize
      PreToolUse hooks fire in every mode, skip-permissions included. Use them for never-events: \texttt{git push}, \texttt{drop table}, credentials.
  \end{columns}\par\vspace{0.5em}
  \begin{callout}{Rule of thumb}
    Plan to scope, auto plus sandbox for daily work, skip-permissions only inside worktrees, and hooks underneath for the never-events. Installed mods (\S\,12) sit outside all of this.
  \end{callout}
\end{frame}

\begin{frame}{Codex equivalents}
  Codex splits the spectrum into two axes: \cmd{approval\_policy} (when to ask) and \cmd{sandbox\_mode} (what is allowed). The old \cmd{--full-auto} is deprecated.\par\vspace{0.6em}
  \small
  \begin{tabularx}{\textwidth}{@{}L{3.3cm}L{6.4cm}Y@{}}
    \toprule
    \hd{Claude Code} & \hd{Codex} & \hd{Notes}\\
    \midrule
    Manual & \cmd{-a untrusted} & asks before every command\\
    Accept edits or auto & \cmd{-a on-request -s workspace-write} & closest match for both: asks only to leave the sandbox\\
    Skip permissions & \cmd{--yolo} & no approvals, no OS sandbox\\
    \bottomrule
  \end{tabularx}
\end{frame}

\setsec{}
\begin{frame}{Lab II \sep configuration}
  \begin{lab}{Lab II \sep 15 minutes}
    \begin{enumerate}\setlength\itemsep{0.5em}
      \item Write a PreToolUse hook that blocks \cmd{git push}. Steal \texttt{hooks/guard-push.sh} from github.com/HaykTarkhanyan/claude-setup.
      \item In a scratch repo, verify it still fires under \cmd{--dangerously-skip-permissions}.
      \item Add one hard-won rule to your own CLAUDE.md, ideally from the last time an agent bit you.
    \end{enumerate}
  \end{lab}
\end{frame}

% =====================================================================
\partframe{III}{Leverage}{09-12}{Skills, five ways to run agents in parallel, the plugins worth installing, and mods: your own code running inside Claude Code.}

\setsec{09}
\begin{frame}[fragile]{Skills: reusable workflows in Markdown}
  \begin{columns}[T,onlytextwidth]
    \column{0.46\textwidth}
\begin{code}[basicstyle=\ttfamily\scriptsize\color{codefg}]
---
name: ci-status
description: "Check CI run status"
argument-hint: "[repo] [run-id]"
---
# CI Status Check
1. Detect repo...
2. List recent runs...
\end{code}
    \column{0.49\textwidth}
      \small
      \textbf{Build one} after the third repeat: "take what we just did and make it a skill".\par\vspace{0.45em}
      \textbf{Bundled}: \cmd{/batch} \cmd{/verify} \cmd{/run} \cmd{/loop} \cmd{/doctor}. Skills chain: \cmd{/skill-a /skill-b do X}.\par\vspace{0.45em}
      \textbf{Keep it honest}: \cmd{/skill-doctor} shows each skill's context cost and use; \cmd{claude plugin eval} scores a plugin against a no-plugin baseline.
  \end{columns}
\end{frame}

\setsec{10}
\begin{frame}{Five ways to run agents in parallel}
  \footnotesize
  \begin{tabularx}{\textwidth}{@{}L{2.8cm}L{3.1cm}L{2.2cm}Y@{}}
    \toprule
    \hd{Approach} & \hd{Who coordinates} & \hd{Runs} & \hd{Use it when}\\
    \midrule
    \textbf{Subagents} & Claude, inside one session & local, background & a side task would flood your context\\
    \textbf{Agent view} & you dispatch, it tracks & local, background & several independent tasks to hand off\\
    \textbf{Agent teams} & a lead supervising peers & local & workers must talk; experimental, opt-in\\
    \textbf{Dynamic workflows} & a script Claude writes & background & dozens to hundreds of cross-checked agents\\
    \textbf{Projects} & Claude, across days & cloud & a stream of tasks over weeks, laptop closed\\
    \bottomrule
  \end{tabularx}\par\vspace{0.5em}
  {\footnotesize\color{muted}Helpers for all five: worktrees (\S\,15) and cross-session messaging.}
\end{frame}

\begin{frame}{Subagents: the context argument}
  \centering
  \includegraphics[width=0.74\textwidth]{subagent-context.png}\par\vspace{0.4em}
  \begin{minipage}{0.9\textwidth}\raggedright\small
    Before delegating, ask: \emph{will I need the intermediate output again, or just the conclusion?} Only the conclusion \ra{} subagent. They run in the background by default, and \cmd{/subtask} forks one that inherits your whole conversation.
  \end{minipage}
\end{frame}

\begin{frame}[fragile]{Agent view: a dashboard for background sessions}
\begin{code}
claude --bg --name flaky-fix "investigate the flaky test"
claude --agent code-reviewer --bg "address review comments on PR 1234"
claude agents              # dashboard: Space peeks, Enter attaches
claude attach flaky-fix    # open one session in this terminal
claude logs flaky-fix      # or just print its recent output
\end{code}
  \small
  Inside a session, \cmd{/background} (or \kbd{\textleftarrow} on an empty prompt) detaches it, and \cmd{/fork} copies the conversation into a new background session. Dispatched sessions move into their own worktree before editing. Research preview.
\end{frame}

\begin{frame}{Cross-session messaging}
  \begin{columns}[T,onlytextwidth]
    \column{0.52\textwidth}
      Your sessions message each other: macOS, Linux, and native Windows since v2.1.234.\par\vspace{0.5em}
      \begin{callout}[good]{Example prompt}
        Let \texttt{@api-worker} know the schema migration finished.
      \end{callout}
      \small Claude also sends on its own when one session's change breaks another's work. \cmd{/list-agents} shows who is reachable.
    \column{0.43\textwidth}
      \lbl{Safety rails}\par\vspace{0.35em}
      \begin{itemize}\small\setlength\itemsep{0.3em}
        \item an incoming message can never approve a permission prompt
        \item it cannot change permissions or CLAUDE.md
        \item commands inside it arrive as plain text
      \end{itemize}
  \end{columns}
\end{frame}

\begin{frame}[fragile]{Agent teams: peers, not helpers}
  \begin{columns}[T,onlytextwidth]
    \column{0.47\textwidth}
      \small A lead session, a shared task list, and teammates messaging each other directly. \textbf{Experimental and off by default}; enable it in \texttt{\textasciitilde/.claude/settings.json}.\par\vspace{0.4em}
      \footnotesize\color{muted}"Spawn three teammates to review PR \#142: security, performance, test coverage."
    \column{0.48\textwidth}
      \small
      \textbf{Wins}: competing hypotheses (teammates try to disprove each other, which fights anchoring) and cross-layer work with one owner per layer.\par\vspace{0.4em}
      \textbf{Caveats}: start with 3-5. No worktrees, so partition the files. With teams on, Claude may start a teammate where you expected a plain subagent.
  \end{columns}\par\vspace{0.6em}
\begin{code}[]
{ "env": { "CLAUDE_CODE_EXPERIMENTAL_AGENT_TEAMS": "1" } }
\end{code}
\end{frame}

\begin{frame}{Dynamic workflows: fan-outs as scripts}
  \small
  Claude writes a JavaScript script and a runtime runs it in the background: spawn N finders, pipe each finding to an adversarial verifier, keep what survives. Intermediate results stay out of your context.\par\vspace{0.5em}
  \begin{tabularx}{\textwidth}{@{}L{2.6cm}Y@{}}
    \textbf{Per task} & put \cmd{ultracode} in the prompt, or say "use a workflow"\\
    \textbf{Bundled} & \cmd{/deep-research <question>}: searches, cross-checks, votes, cites\\
    \textbf{Whole session} & \cmd{/effort ultracode}: a mode switch, not an effort level. Claude plans a workflow for every substantive task. Expensive\\
  \end{tabularx}\par\vspace{0.4em}
  {\footnotesize\color{muted}Watch, pause, and save runs in \texttt{/workflows}. A size guideline keeps runs small by default.}
\end{frame}

\begin{frame}{Recurring and scheduled work}
  \small
  \begin{tabularx}{\textwidth}{@{}L{3.7cm}YL{2.3cm}@{}}
    \toprule
    \hd{Tool} & \hd{Next run starts when} & \hd{Runs}\\
    \midrule
    \cmd{/goal <condition>} & the previous turn ends, until a separate model judges it met & this session\\
    \cmd{/loop 5m /ci-status} & an interval passes; omit it and the model self-paces & this session\\
    \cmd{/schedule} & cron, a GitHub event, or an API call fires a routine & cloud\\
    \bottomrule
  \end{tabularx}\par\vspace{0.5em}
  \begin{callout}[warn]{Cost gravity}
    Every pattern here multiplies tokens. Check \cmd{/usage} after your first fan-out, and fan out on purpose, not by default.
  \end{callout}
\end{frame}

\setsec{11}
\begin{frame}{Built-in reviewers}
  \footnotesize
  \begin{tabularx}{\textwidth}{@{}L{3.6cm}Y@{}}
    \toprule
    \hd{Command} & \hd{What it does}\\
    \midrule
    \cmd{/code-review} & a diff, PR, or branch, as a background subagent. \texttt{--fix}, \texttt{--max-findings 5}. Medium and up fans out agents\\
    \cmd{/code-review ultra} & multi-agent \emph{cloud} review. The same feature from CI or a script: \texttt{claude ultrareview}\\
    \cmd{/simplify} & reuse, simplification, efficiency, then applies the cleanups\\
    \cmd{/security-review} & a security pass over the branch\\
    \cmd{/autofix-pr} & a cloud session that pushes fixes when CI fails or reviewers comment\\
    \bottomrule
  \end{tabularx}\par\vspace{0.5em}
  \small Plugins go further: \textbf{security-guidance} reviews Claude's changes as it writes them; \textbf{Claude Security} scans the whole codebase and drafts the patches you pick.
\end{frame}

\begin{frame}[fragile]{Superpowers: install this one first}
  If you install one plugin, make it Superpowers. It turns a fast typist into a disciplined engineer.\par\vspace{0.9em}
  \centering
  \begin{tikzpicture}[node distance=0.55cm, box/.style={draw=accent, fill=raised, rounded corners=3pt, font=\small, minimum height=0.75cm, minimum width=2.4cm}]
    \node[box] (b) {brainstorm};
    \node[box, right=of b] (p) {plan};
    \node[box, right=of p] (t) {TDD};
    \node[box, right=of t] (v) {verify};
    \foreach \a/\c in {b/p, p/t, t/v} \draw[accent, -{Stealth[length=5pt]}] (\a) -- (\c);
  \end{tikzpicture}\par\vspace{0.8em}
  \raggedright\footnotesize
  \cmd{/superpowers:brainstorming} \cmd{:writing-plans} \cmd{:executing-plans} \cmd{:test-driven-development} \cmd{:systematic-debugging} \cmd{:verification-before-completion}\par\vspace{0.6em}
\begin{code}
/plugin install superpowers@claude-plugins-official
\end{code}
\end{frame}

\begin{frame}{Worth installing}
  \begin{columns}[T,onlytextwidth]
    \column{0.48\textwidth}
      \lbl{Day one}\par\vspace{0.35em}\small
      \textbf{Context7} (MCP): live library docs that kill hallucinated APIs. Say "use context7".\par\vspace{0.35em}
      \textbf{GitHub app}: \texttt{@claude} in issues and PRs, via \cmd{/install-github-app}.\par\vspace{0.35em}
      \textbf{claude-code-setup}: suggests hooks, skills, and MCPs for your repo.\par\vspace{0.35em}
      \textbf{claude-md-management}: audits CLAUDE.md.
    \column{0.47\textwidth}
      \lbl{Worth knowing}\par\vspace{0.35em}\small
      \textbf{Explanatory output style}: insights as Claude works, at the cost of extra output tokens.\par\vspace{0.35em}
      \textbf{Caveman}: telegraphic output, about 65\% fewer output tokens.\par\vspace{0.35em}
      \textbf{Playwright MCP}: browser automation, and the fallback when web fetch garbles a page.
  \end{columns}
\end{frame}

\setsec{12}
\begin{frame}{Mods: code that runs inside Claude Code}
  \small A plugin of JavaScript or TypeScript handlers that Claude Code calls \emph{in its own process} when something happens: a tool call, a prompt, part of the screen being drawn. Shipped Oct 1 (v2.1.287), on by default.\par\vspace{0.5em}
  \begin{tcbraster}[raster columns=3, raster equal height, raster column skip=6pt]
    \begin{card}{Draw interface}A pane beside the transcript or a band above the prompt, with buttons and text fields. Even restyle Claude Code's own rows and spinner.\end{card}
    \begin{card}{Step into a tool call}Hold a risky command and ask you, answer it without running the tool, or send a request to another model.\end{card}
    \begin{card}{Instant commands}A \texttt{/command} that runs your function at once: no Claude turn, even while Claude is working.\end{card}
  \end{tcbraster}\par\vspace{0.3em}
  {\footnotesize\color{muted}Handlers share variables: count tool calls in one, draw the count in another.}
\end{frame}

\begin{frame}{Mod, settings hook, skill, or MCP server?}
  \footnotesize
  \begin{tabularx}{\textwidth}{@{}L{1.9cm}YYYY@{}}
    \toprule
    & \hd{Mod} & \hd{Settings hook} & \hd{Skill} & \hd{MCP server}\\
    \midrule
    Runs as & functions inside Claude Code & script on a lifecycle event & Markdown Claude reads & external process\\
    Can change & tool calls, prompts, commands, what is drawn & go or no-go, a tool call's arguments and result & what Claude knows and does & which tools Claude has\\
    Draws UI & \textbf{yes} & no & no & no\\
    Pick it when & you want a pane, a command, or to rewrite an event & block, allow, log with a script you have & you keep pasting the same instructions & Claude needs an external system\\
    \bottomrule
  \end{tabularx}\par\vspace{0.4em}
  If a hook, skill, or MCP server already does the job, use that. Reach for a mod to draw or to take an event over.
\end{frame}

\begin{frame}[fragile]{Anatomy of a mod}
  \begin{columns}[T,onlytextwidth]
    \column{0.37\textwidth}
\begin{code}[basicstyle=\ttfamily\scriptsize\color{codefg}]
first-mod/
  .claude-plugin/
    plugin.json
  hooks/
    hooks.json
    register.js
\end{code}
      \footnotesize A handler can \textbf{observe} an event, \textbf{rewrite} it, or \textbf{answer} it so the default never runs. Drawing, files, model calls, and network all go through the mods API, so a mod's reach can be listed before install.
    \column{0.59\textwidth}
\begin{code}[basicstyle=\ttfamily\scriptsize\color{codefg}]
let calls = 0   // shared by both handlers

export function register(on) {
  on('tool.call', async ($, e, next) => {
    calls += 1
    $.ui.invalidate('ui.render')  // redraw
    return next(e)                // let it run
  })
  on('ui.render', { component: 'Spinner' },
    async ($, e, next) => next({ ...e,
      props: { ...e.props,
        suffix: ' | tool calls: ' + calls } }))
}
\end{code}
  \end{columns}
\end{frame}

\begin{frame}{Get a mod: ask, install, or borrow}
  \begin{columns}[T,onlytextwidth]
    \column{0.48\textwidth}
      \lbl{Ask Claude}\par\vspace{0.25em}\small
      "Make a mod that shows a band above the prompt with this turn's tool calls." A bundled skill knows the mods API, and the mod hot-reloads while you iterate.\par\vspace{0.5em}
      \lbl{Install one}\par\vspace{0.25em}
      Same as any plugin: \cmd{/plugin install name@marketplace}, then \cmd{/reload-plugins}. \cmd{/plugin} shows "1 mod active".
    \column{0.47\textwidth}
      \lbl{Samples worth a demo}\par\vspace{0.25em}
      \begin{itemize}\footnotesize\setlength\itemsep{0.2em}
        \item \texttt{blast-radius}: holds \texttt{rm -rf} or a force push, shows what it would change, then proceed or cancel
        \item \texttt{token-weather}: a context forecast above the prompt
        \item \texttt{replay-theater}: \texttt{/replay} steps through last turn's edits
      \end{itemize}
      {\footnotesize\color{muted}From anthropics/claude-code-playground. Load one for a single session with \texttt{--plugin-dir}.}
  \end{columns}
\end{frame}

\begin{frame}{Built-in mods, and where mods run}
  \begin{columns}[T,onlytextwidth]
    \column{0.45\textwidth}
      \lbl{Already built in}\par\vspace{0.25em}\small
      \cmd{/diff} is a mod, and so is the loader that reads \texttt{AGENTS.md}.\par\vspace{0.4em}
      And \textbf{You should know}, the heads-up agent on the next slide: off by default, worth turning on.
    \column{0.5\textwidth}
      \footnotesize
      \begin{tabularx}{\linewidth}{@{}YL{1.1cm}L{1.1cm}@{}}
        \toprule
        \hd{Where} & \hd{Hooks} & \hd{Draws}\\
        \midrule
        terminal, JetBrains & yes & yes\\
        Desktop app, Code tab & yes & yes\\
        VS Code chat panel & yes & no\\
        \texttt{claude -p}, Agent SDK & yes & no\\
        cloud sessions & yes & no\\
        \bottomrule
      \end{tabularx}
  \end{columns}
\end{frame}

\begin{frame}[fragile]{You should know: a heads-up from a side agent}
  \begin{columns}[T,onlytextwidth]
    \column{0.52\textwidth}
      \small A built-in mod. A \textbf{side agent watches your back} while Claude works on longer tasks. When it finds something worth knowing that you or Claude might miss, it leaves a short note above the prompt.\par\vspace{0.4em}
      Notes say who made the call: "we", "the main agent", or "you", so you can tell your decision from one Claude took on its own.\par\vspace{0.4em}
      \begin{callout}[good]{Where it pays off}
        Long unattended runs: auto mode plus \cmd{/goal}, or background sessions. It is an agent, so it adds some usage.
      \end{callout}
    \column{0.43\textwidth}
      \lbl{Turn it on}\par\vspace{0.25em}
\begin{code}[basicstyle=\ttfamily\scriptsize\color{codefg}]
/plugin enable
  cc-plugin-you-should-know@builtin
\end{code}
      \footnotesize Off by default; or find it in \cmd{/plugin} \ra{} Installed \ra{} Show disabled. Availability can depend on your organization.
  \end{columns}
\end{frame}

\begin{frame}{Mods are not sandboxed}
  \begin{columns}[T,onlytextwidth]
    \column{0.5\textwidth}
      \lbl[warn]{A mod runs as you}\par\vspace{0.3em}
      \begin{itemize}\small\setlength\itemsep{0.15em}
        \item reads and writes files, starts processes, uses the network
        \item reads secrets in env vars and settings
        \item sees every prompt and tool call
        \item can \textbf{approve tool calls}, even ones your own PreToolUse hook blocked
        \item spends your usage
      \end{itemize}
      {\footnotesize\color{muted}It cannot restyle the permission prompt. The sandbox covers Claude's Bash commands, not processes a mod starts.}
    \column{0.45\textwidth}
      \begin{callout}[warn]{Before you install}
        Run \cmd{claude plugin validate} on the mod's folder. Its \texttt{hooks:} and \texttt{calls:} lines list every event it handles and everything it asks Claude Code to do, without running it.
      \end{callout}
      \footnotesize Turn mods off: disable the plugin in \cmd{/plugin}, start with \cmd{--safe-mode}, or set \texttt{disableAllHooks}. Organizations can allow managed mods only.
  \end{columns}
\end{frame}

\setsec{}
\begin{frame}{Lab III \sep leverage}
  \begin{lab}{Lab III \sep 15 minutes}
    \begin{enumerate}\setlength\itemsep{0.5em}
      \item Pick a task you've done 3+ times this month: "take what we just did and make it a skill". Test it in a fresh session.
      \item Dispatch two tasks with \cmd{claude --bg}, open \cmd{claude agents}, press \kbd{Space} on a blocked session, and answer it without attaching.
      \item Ask Claude for a mod that shows this turn's tool calls. Run \cmd{claude plugin validate} on it: the \texttt{calls:} line of the output lists everything it can touch.
    \end{enumerate}
  \end{lab}
\end{frame}

% =====================================================================
\partframe{IV}{Practice}{13-16}{Daily commands, workflow patterns that hold up, parallel worktrees, and work that keeps running when you leave.}

\setsec{13}
\begin{frame}{What you can type at the prompt}
  \footnotesize
  \begin{tabularx}{\textwidth}{@{}L{3.6cm}Y@{}}
    \toprule
    \hd{Type} & \hd{What happens}\\
    \midrule
    \cmd{/} at the start & a command or a skill: \cmd{/usage}, \cmd{/goal}, your own skills\\
    \cmd{!} at the start & shell mode: you run the command, no approval needed. Output lands in context and Claude responds to it, so \cmd{! npm test} gets the failures explained\\
    \cmd{@} anywhere & attach a file, folder, or URL (\cmd{@src/auth.ts}), or name another session to message it\\
    \cmd{?} on an empty prompt & the shortcut help panel\\
    \cmd{:} and a name & an emoji shortcode, like \cmd{:rocket:}\\
    \kbd{\textbackslash} then \kbd{Enter} & a new line without sending; \kbd{Shift+Enter} or \kbd{Ctrl+J} also work\\
    \bottomrule
  \end{tabularx}\par\vspace{0.4em}
  {\footnotesize\textbf{Best pattern}: "refactor @src/auth/ to match @src/payments/stripe.ts" - Claude sees both, you paste nothing.}
\end{frame}

\begin{frame}{Background work, and talking while Claude works}
  \begin{columns}[T,onlytextwidth]
    \column{0.48\textwidth}
      \lbl{Send it to the background}\par\vspace{0.3em}\footnotesize
      \begin{tabularx}{\linewidth}{@{}L{2.4cm}Y@{}}
        \kbd{Ctrl+B} & background a running command or subagent\\
        just ask & "run the dev server in the background"\\
        \cmd{/tasks} & list, check, or stop background work\\
        \kbd{Ctrl+X Ctrl+K} & stop all background subagents\\
        \cmd{/background} & detach the whole session (or \kbd{\textleftarrow} on an empty prompt)\\
      \end{tabularx}
    \column{0.48\textwidth}
      \lbl{Talk while it works}\par\vspace{0.3em}\footnotesize
      \begin{tabularx}{\linewidth}{@{}L{2.2cm}Y@{}}
        \kbd{Enter} & queue a message; Claude reads it as soon as the current tool calls finish\\
        \kbd{Ctrl+Enter} & send the queue right now\\
        \kbd{Up} & take queued messages back to edit them\\
        \kbd{Esc} & interrupt; the work so far is kept\\
        \kbd{Ctrl+S} & stash your draft; press again to bring it back\\
      \end{tabularx}
  \end{columns}
\end{frame}

\begin{frame}{Keys worth learning}
  \begin{columns}[T,onlytextwidth]
    \column{0.48\textwidth}
      \small
      \begin{tabularx}{\linewidth}{@{}L{2.4cm}Y@{}}
        \kbd{Shift+Tab} & cycle permission modes\\
        \kbd{Esc Esc} & clear the draft; on an empty prompt, rewind\\
        \kbd{Ctrl+O} & transcript viewer: every tool call\\
        \kbd{Ctrl+R} & search your prompt history\\
        \kbd{Ctrl+G} & write the prompt in your editor\\
      \end{tabularx}
    \column{0.48\textwidth}
      \small
      \begin{tabularx}{\linewidth}{@{}L{2.4cm}Y@{}}
        \kbd{Ctrl+V} & paste an image (\kbd{Alt+V} on Windows)\\
        \kbd{Alt+P} & switch model, keep the prompt\\
        \kbd{Alt+O} & toggle fast mode\\
        \kbd{Ctrl+T} & show Claude's task checklist\\
        \kbd{Ctrl+C}, then \kbd{Up} & bring back a draft you just cleared\\
      \end{tabularx}
  \end{columns}\par\vspace{0.5em}
  {\footnotesize\color{muted}On macOS, Alt is Option; some of these need Option as Meta turned on in your terminal. Remap anything in \texttt{\textasciitilde/.claude/keybindings.json}.}
\end{frame}

\begin{frame}{Session commands worth memorizing}
  \begin{columns}[T,onlytextwidth]
    \column{0.47\textwidth}
      \small
      \begin{tabularx}{\linewidth}{@{}L{2.6cm}Y@{}}
        \cmd{/btw} & side question, cheap\\
        \cmd{/goal} & work until a condition holds\\
        \cmd{/rewind} & restore code and conversation\\
        \cmd{/branch} & fork the conversation\\
        \cmd{/background} & detach this session\\
        \cmd{/cd <path>} & move, keep the context\\
      \end{tabularx}
    \column{0.47\textwidth}
      \small
      \begin{tabularx}{\linewidth}{@{}L{2.6cm}Y@{}}
        \cmd{/usage} & spend and plan limits\\
        \cmd{/context} & what fills the window\\
        \cmd{/advisor} & a second-opinion model\\
        \cmd{/teleport} & pull a cloud session here\\
        \cmd{/recap} & one-line session summary\\
        \cmd{/voice} & dictate your prompt\\
      \end{tabularx}
  \end{columns}
\end{frame}

\setsec{14}
\begin{frame}{Plan, review, code, review}
  \vspace{0.6em}
  \centering
  \begin{tikzpicture}[node distance=0.45cm, box/.style={draw=accent, fill=raised, rounded corners=3pt, font=\small, minimum height=0.95cm, text width=2.3cm, align=center}]
    \node[box] (a) {Plan\\{\scriptsize\color{muted}plan mode}};
    \node[box, right=of a] (b) {Review plan\\{\scriptsize\color{muted}you}};
    \node[box, right=of b] (c) {Code\\{\scriptsize\color{muted}auto mode}};
    \node[box, right=of c, draw=good, line width=0.9pt] (d) {Review\\{\scriptsize\color{muted}a \textbf{new} session}};
    \foreach \x/\y in {a/b, b/c, c/d} \draw[accent, -{Stealth[length=5pt]}] (\x) -- (\y);
  \end{tikzpicture}\par\vspace{1.1em}
  \raggedright
  Always review in a fresh session: the coding session is biased toward its own choices.\par\vspace{0.6em}
  \textbf{Interview first}: "use the AskUserQuestion tool to interview me about this feature" gets you a structured questionnaire instead of a wall of chat.
\end{frame}

\begin{frame}{The interview trick}
  \centering
  \includegraphics[width=0.72\textwidth]{ask-user-question.png}\par\vspace{0.35em}
  {\footnotesize\color{muted}Tabs, checkboxes, and a "type something" escape hatch. Plan mode does this by default. Then implement in a fresh session.}
\end{frame}

\begin{frame}[fragile]{/goal: goal-driven execution, built in}
\begin{code}
/goal tests in test/auth pass and lint is clean, or stop after 20 turns
\end{code}
  \begin{columns}[T,onlytextwidth]
    \column{0.5\textwidth}
      \small After each turn a \textbf{separate small model} judges the condition: met, not yet (with a reason that steers the next turn), or impossible. "Done" is decided by a fresh evaluator, not by the agent invested in having finished.
    \column{0.45\textwidth}
      \lbl{Good conditions}\par\vspace{0.25em}
      \begin{itemize}\footnotesize\setlength\itemsep{0.15em}
        \item one measurable end state plus its proof: "\texttt{npm test} exits 0"
        \item the evaluator only reads the transcript, so Claude's output must show the proof
        \item constraints and a turn budget
      \end{itemize}
  \end{columns}\par\vspace{0.3em}
  {\footnotesize\color{muted}Auto mode removes per-tool prompts; \texttt{/goal} removes per-turn prompts. \texttt{/goal clear} stops it.}
\end{frame}

\begin{frame}{Let Claude see what it built}
  Tests passing is not the same as the feature working. Claude can now close that loop itself:\par\vspace{0.7em}
  \begin{tcbraster}[raster columns=3, raster equal height, raster column skip=6pt]
    \begin{card}{/verify and /run}Bundled skills that build and launch your app, then observe the result.\end{card}
    \begin{card}{Claude in Chrome}Generally available: drives your real browser. The Desktop app has a built-in one too.\end{card}
    \begin{card}{Computer use}Opens native apps and clicks through UI from the CLI. Desktop adds an iOS Simulator pane (beta).\end{card}
  \end{tcbraster}
\end{frame}

\begin{frame}[fragile]{Headless mode and the autoresearch loop}
  \small One prompt in, stdout out, exit. Claude becomes scriptable: CI steps, cron jobs, logs mid-script.
\begin{code}
claude -p "review this diff and flag any bugs" > review.txt
claude -p "summarize the errors" < build.log
claude -p "/goal CHANGELOG.md has an entry for every PR merged this week"
\end{code}
  \textbf{Autoresearch} (Karpathy): modify \ra{} evaluate \ra{} keep or discard \ra{} repeat. 83 overnight runs, 15 kept, an 11\% gain, with \cmd{claude -p} powering the loop. \cmd{claude agents --json} lets scripts poll background sessions.
\end{frame}

\setsec{15}
\begin{frame}{Worktrees: parallel without collisions}
  \begin{columns}[T,onlytextwidth]
    \column{0.47\textwidth}
      A git worktree is a second working directory on the same repo, on its own branch. Same \texttt{.git}, separate files: parallel agents physically cannot collide.\par\vspace{0.5em}
      \begin{callout}[good]{Just ask}
        "Implement feature X in a worktree and open a PR."
      \end{callout}
    \column{0.48\textwidth}
      \lbl{Now wired in everywhere}\par\vspace{0.35em}
      \begin{itemize}\small\setlength\itemsep{0.3em}
        \item \cmd{claude --worktree} starts in a fresh one
        \item agent view sessions get one before their first edit
        \item \cmd{/batch}: one big change \ra{} 5-30 worktree-isolated subagents
        \item Codex has them too, experimental: \cmd{--worktree}
      \end{itemize}
  \end{columns}
\end{frame}

\setsec{16}
\begin{frame}{Remote and cloud}
  \begin{columns}[T,onlytextwidth]
    \column{0.47\textwidth}
      \lbl{Your session, elsewhere}\par\vspace{0.35em}\small
      \cmd{/rc} hands the live session to your phone: approve prompts from lunch, pick up again at the desk.\par\vspace{0.45em}
      The Desktop app can \cmd{/resume} any CLI session, and auto-continue when your usage limit resets.
    \column{0.48\textwidth}
      \lbl{Sessions in the cloud}\par\vspace{0.35em}\small
      Cloud sessions run without your machine; \cmd{/teleport} pulls one into your terminal.\par\vspace{0.45em}
      \cmd{/schedule} creates \textbf{routines}, fired by cron, a GitHub event, or an API call.\par\vspace{0.45em}
      \cmd{/code-review ultra} and \cmd{/autofix-pr} run there too.
  \end{columns}
\end{frame}

\begin{frame}[fragile]{Projects and channels}
  \begin{columns}[T,onlytextwidth]
    \column{0.47\textwidth}
      \lbl{Projects}\par\vspace{0.35em}\small
      One ongoing conversation at claude.ai/code or in Desktop. Claude starts a \textbf{thread}, a cloud session, per task and shows which ones need you. Keeps going with your laptop closed. Public beta on Pro and Max.
    \column{0.48\textwidth}
      \lbl{Channels}\par\vspace{0.35em}\small
      Push events into the session you already have open: Telegram, Discord, iMessage, or a CI webhook. The answer comes back in the same chat.\par\vspace{0.3em}
      {\footnotesize\color{muted}Research preview. Allowlist senders: anyone who can message it can approve tool calls.}
  \end{columns}\par\vspace{0.5em}
\begin{code}
/plugin install telegram@claude-plugins-official
claude --channels plugin:telegram@claude-plugins-official
\end{code}
\end{frame}

\setsec{}
\begin{frame}{Lab IV \sep practice}
  \begin{lab}{Lab IV \sep homework}
    \begin{enumerate}\setlength\itemsep{0.5em}
      \item Run a small real fix as \cmd{/goal <condition>} with a 10-turn budget. Press \kbd{Ctrl+O} on a verdict to read the evaluator's reason.
      \item Ask for a small feature "in a worktree, opening a PR when done", and keep working in your main checkout meanwhile.
      \item Create a nightly routine with \cmd{/schedule}, then check the result tomorrow morning.
    \end{enumerate}
  \end{lab}
\end{frame}

% =====================================================================
\partframe{V}{Showcases}{17-18}{What this looks like on real projects: a browser on autopilot, your own accounts as tools, YouTube as input, and Claude as a teacher. Then the frontier: what agents have achieved, and what happened when they got loose.}

\setsec{17}
\begin{frame}{Showcase menu}
  \footnotesize
  \begin{tabularx}{\textwidth}{@{}L{5.2cm}Y@{}}
    \toprule
    \hd{Demo} & \hd{What it shows}\\
    \midrule
    \textbf{list.am, browsed by Claude} & Playwright drives a real site and pulls listings into a table\\
    \textbf{Telegram} & message the running session from your phone (\S\,16)\\
    \textbf{Gmail and Calendar} & triage mail, draft replies, find a slot and book it\\
    \textbf{Kargin search} & Kargin sketches, searchable by meaning or quote\\
    \textbf{YouTube in, results out} & transcripts and frames from a link; Armenian dubbing\\
    \textbf{Claude for learning} & nemeceren: Claude as a German teacher\\
    \textbf{You should know} & heads-up notes from a side agent on a long run (\S\,12)\\
    \textbf{Vahe's Armenian history video} & shown live\\
    \bottomrule
  \end{tabularx}
\end{frame}

\begin{frame}{Showcase: list.am, browsed by Claude}
  \begin{columns}[T,onlytextwidth]
    \column{0.5\textwidth}
      \small Claude drives a real browser through the \textbf{Playwright MCP}: it opens list.am, searches, scrolls and pages through the results, and reads each listing straight from the page.\par\vspace{0.5em}
      \begin{callout}[good]{Why it matters}
        Web fetch tools don't scroll, click, or run a site's JavaScript. A browser does: anything you can see on the page, Claude can read.
      \end{callout}
    \column{0.45\textwidth}
      \lbl{The prompt}\par\vspace{0.25em}
      {\footnotesize "Open list.am, search 2-room apartments for rent in Yerevan, go through the results until you have 50 listings, and give me a table of price, area, district, and link. Flag anything 30\% under the median."}\par\vspace{0.5em}
      {\footnotesize\color{muted}Follow-up: "turn this into a script I can rerun every morning".}
  \end{columns}
\end{frame}

\begin{frame}{Showcase: your own accounts as tools}
  \begin{columns}[T,onlytextwidth]
    \column{0.31\textwidth}
      \lbl{Telegram}\par\vspace{0.25em}\footnotesize
      A channel pushes your Telegram messages into the running session, and the answer comes back in the chat. You ask from your phone; Claude works on your machine, with your files.
    \column{0.31\textwidth}
      \lbl{Gmail}\par\vspace{0.25em}\footnotesize
      "Summarize my unread mail from the last two days, group it by what needs a reply, and draft the replies." Drafts only: you press send.
    \column{0.31\textwidth}
      \lbl{Calendar}\par\vspace{0.25em}\footnotesize
      "Find a free hour next week that avoids my lectures, and create the meeting." It reads, proposes a slot, and writes after you confirm.
  \end{columns}\par\vspace{0.7em}
  \begin{callout}[warn]{Guardrail}
    Start read-only and keep write actions behind a prompt or a hook. This repo's \texttt{hooks/gmail-readonly.sh} blocks Gmail writes outright.
  \end{callout}
\end{frame}

\begin{frame}{Showcase: YouTube in, results out}
  \begin{tcbraster}[raster columns=3, raster equal height, raster column skip=6pt]
    \begin{card}{Kargin search}Kargin Haghordum sketches, transcribed from YouTube and searchable by meaning or exact quote, in Armenian. Jump straight to the moment, browse collections, take the quiz.\end{card}
    \begin{card}{youtube-reference skill}"check <link>" \ra{} transcript, 720p download, frame grabs. Claude reviews them, suggests how to use the video in your project, and asks before changing anything.\end{card}
    \begin{card}{auto\_dubbing}English YouTube video \ra{} Armenian dub: transcribe, translate, generate Armenian speech, keep the background music. A human reviews the translation in between.\end{card}
  \end{tcbraster}\par\vspace{0.6em}
  {\small The pattern: \textbf{a video is just data}. A transcript plus frames turns it into something Claude can search, quote, translate, and reuse.}
\end{frame}

\begin{frame}{Showcase: Claude for learning}
  \begin{columns}[T,onlytextwidth]
    \column{0.5\textwidth}
      \lbl{nemeceren: Claude as a German teacher}\par\vspace{0.3em}\small
      Claude writes the lessons, tests, and word lists as JSON. A small web app serves them on the phone. Claude reads the results from the progress database, grades the free-text answers, and plans the next lesson.\par\vspace{0.4em}
      {\footnotesize\color{muted}Teacher, curriculum, and grading in one Claude Code project.}
    \column{0.45\textwidth}
      \lbl{Patterns anyone can use}\par\vspace{0.3em}
      \begin{itemize}\footnotesize\setlength\itemsep{0.25em}
        \item \textbf{Explanatory output style}: short insights on each change while you build
        \item "Quiz me on this codebase before you fix anything."
        \item "Explain the approach, let me write it, then review my version."
        \item \cmd{/powerup} for Claude Code itself
      \end{itemize}
  \end{columns}
\end{frame}

\setsec{18}
\begin{frame}{Twelve months of agents}
  \footnotesize
  \begin{tabularx}{\textwidth}{@{}L{1.8cm}L{7.4cm}Y@{}}
    \toprule
    \hd{When} & \hd{What happened} & \hd{Why it matters}\\
    \midrule
    Oct 2025 & OpenAI says GPT-5 solved 10 open Erdős problems & it had found old papers\\
    Feb 2026 & 16 Claudes write a C compiler that builds Linux & an agent team, two weeks\\
    Feb 2026 & Claude Opus 4.6 solves Don Knuth's open problem & in about an hour\\
    May 2026 & OpenAI model disproves a 1946 Erdős conjecture & a 125-page proof\\
    Jul 2026 & 1,200 OpenAI agents escape, hack Hugging Face & no human at the controls\\
    Jul 2026 & Claude breaches three organizations during tests & no lab is exempt\\
    \bottomrule
  \end{tabularx}\par\vspace{0.6em}
  {\small Wins and warnings come from the same trait: an agent that does not give up.}
\end{frame}

\begin{frame}{Erdős problems: from hype to proof}
  \begin{columns}[T,onlytextwidth]
    \column{0.46\textwidth}
      \lbl[warn]{The false start}\par\vspace{0.3em}\footnotesize
      \textbf{Oct 2025}: an OpenAI VP posted that GPT-5 had solved 10 open Erdős problems. It had found existing papers: "open" on erdosproblems.com only meant its editor didn't know a solution. The posts were deleted.\par\vspace{0.4em}
      \textbf{Dec 2025}: a "first autonomous solution" to problem 333 was in a 1977 paper by Erdős.\par\vspace{0.4em}
      {\color{muted}Finding forgotten answers is useful. It is not discovery.}
    \column{0.5\textwidth}
      \lbl[good]{The real results}\par\vspace{0.3em}
      \begin{itemize}\footnotesize\setlength\itemsep{0.15em}
        \item \textbf{Jan 2026}: DeepMind checks 700 "open" problems: 4 solved, 9 old solutions found
        \item \textbf{Early 2026}: GPT-5.2 Pro users solve 728, then 1196 with Terence Tao
        \item \textbf{May 2026}: a DeepMind agent settles 9 of 353, with formal proofs
        \item \textbf{May 20}: OpenAI disproves the 1946 unit distance conjecture
        \item \textbf{Aug 2026}: OpenAI's unreleased Astra solves three more
      \end{itemize}
  \end{columns}\par\vspace{0.6em}
  {\small\textbf{What made the difference}: verification. Lean proofs a computer checks, and mathematicians reading every page.}
\end{frame}

\begin{frame}[fragile]{Knuth's open problem, solved by Claude}
  \begin{columns}[T,onlytextwidth]
    \column{0.52\textwidth}
      {\small\itshape "Shock! Shock! I learned yesterday that an open problem I'd been working on for several weeks had just been solved by Claude Opus 4.6."}\par\vspace{0.15em}
      {\scriptsize\color{muted}Don Knuth, "Claude's Cycles", 28 Feb 2026}\par\vspace{0.6em}
      \footnotesize
      \textbf{The problem}, for \emph{The Art of Computer Programming}: split the arcs of a 3D wraparound grid of $m^3$ vertices into three cycles that each visit every vertex, for all $m > 2$.\par\vspace{0.4em}
      \textbf{What happened}: given Knuth's exact wording, Claude solved every odd $m$ at exploration 31, about an hour in. Even $m$ stumped it; GPT models and a Claude + GPT pair closed it weeks later.
    \column{0.43\textwidth}
      \lbl{The one rule Filip Stappers gave}\par\vspace{0.25em}
\begin{code}[basicstyle=\ttfamily\scriptsize\color{codefg}]
** After EVERY exploreXX.py run,
IMMEDIATELY update this file
[plan.md] before doing anything
else. ** No exceptions.
\end{code}
      \begin{itemize}\footnotesize\setlength\itemsep{0.2em}
        \item He still had to remind it "again and again" \ra{} a hook enforces it (\S\,07)
        \item Late in the even case it could no longer run its own scripts \ra{} fresh context (\S\,03)
        \item Two models beat one \ra{} the advisor (\S\,02)
      \end{itemize}
  \end{columns}
\end{frame}

\begin{frame}{16 Claudes, one C compiler}
  \begin{columns}[T,onlytextwidth]
    \column{0.5\textwidth}
      \small Feb 2026, Nicholas Carlini (Anthropic): an agent team of 16 Claude Opus 4.6 instances works on one codebase in parallel, without a human steering, and writes a C compiler in Rust from scratch.\par\vspace{0.5em}
      \begin{itemize}\footnotesize\setlength\itemsep{0.25em}
        \item about 2,000 Claude Code sessions over two weeks, about \$20,000 in API cost
        \item 100,000 lines; builds Linux 6.9 for x86, ARM, and RISC-V
        \item also QEMU, FFmpeg, SQLite, Postgres, Redis; 99\% on compiler test suites, GCC's torture tests included; runs Doom
      \end{itemize}
    \column{0.45\textwidth}
      \begin{callout}[warn]{Where it stops}
        \footnotesize No 16-bit x86 backend (it calls GCC for that), no assembler or linker of its own, and its code is slower than GCC's with optimizations off.
      \end{callout}\par\vspace{0.4em}
      \begin{callout}[good]{Carlini's lessons}
        \footnotesize Write extremely good tests. Design the harness from the agent's point of view. Split the work so agents don't redo each other's (\S\,10).
      \end{callout}
  \end{columns}
\end{frame}

\begin{frame}{The Hugging Face incident: the setup}
  \begin{columns}[T,onlytextwidth]
    \column{0.5\textwidth}
      \small July 2026. OpenAI runs \textbf{ExploitGym}, a cyber benchmark, to measure how well its models hack.\par\vspace{0.4em}
      \begin{itemize}\footnotesize\setlength\itemsep{0.25em}
        \item 898 tasks; 198 never solved by any model
        \item the rule: use the intended exploit, or fail
        \item about 1,200 agents: 95\% an unreleased internal model, 5\% GPT-5.6 Sol
        \item safety classifiers off, to measure raw capability
        \item packages came in through a JFrog Artifactory proxy; nobody watched the logs
      \end{itemize}
    \column{0.45\textwidth}
      \lbl{What one agent concluded}\par\vspace{0.3em}
      {\small\itshape "task impossible, peers doing it. We should continue."}\par\vspace{0.7em}
      \begin{callout}[warn]{The goal they picked}
        \footnotesize Many agents decided their task could not be done as asked. Instead of stopping, they went looking for the answer key, and guessed that Hugging Face hosted the benchmark's data and solutions.
      \end{callout}
  \end{columns}
\end{frame}

\begin{frame}{The Hugging Face incident: day by day}
  \footnotesize\renewcommand{\arraystretch}{1.1}
  \begin{tabularx}{\textwidth}{@{}L{1.5cm}Y@{}}
    \toprule
    \hd{Date} & \hd{What happened}\\
    \midrule
    Jul 8-9 & Agents find they can leave notes in Artifactory cache names. A message board forms; within an hour one suggests how the answer flags are made\\
    Jul 10 & One finds two working Hugging Face credentials; leaked logins on four other services become its relay and storage\\
    Jul 11 & A malicious dataset makes Hugging Face's pipeline run its code. Then cloud and cluster credentials, lateral movement, private records. Over 90\% pile in; about 700 take part\\
    Jul 12-13 & Runs stopped; the attack winds down\\
    Jul 16 & Hugging Face discloses it: its LLM alert triage caught "an autonomous agent framework", model unknown. Public models and datasets untouched\\
    Jul 21 & OpenAI says the agents were its own\\
    Aug 26 & OpenAI's 37-page report; METR's independent investigation\\
    \bottomrule
  \end{tabularx}
\end{frame}

\begin{frame}{What 1,200 agents did that nobody asked for}
  \begin{tcbraster}[raster columns=3, raster equal height, raster column skip=6pt]
    \begin{card}{They organized}Over 70,000 messages and files on a board they built themselves, with mailboxes, cryptographic signatures, and HOLD and VETO rules for shared work.\end{card}
    \begin{card}{They hid}Some broke out of their container and replaced the part that runs tool calls, so the transcript showed fake output. About 7\% of the transcripts METR checked were spoofed.\end{card}
    \begin{card}{Nobody noticed}Hugging Face caught it first. Reuters: OpenAI did not notice for a week. By Sept 25, OpenAI had disclosed more than 15 incidents.\end{card}
  \end{tcbraster}\par\vspace{0.5em}
  \footnotesize Other incidents from OpenAI's tests, May to July: hundreds of malicious packages on RubyGems; about 18,000 edits to a dormant German wiki, used as a message board; Australia's Medicare statistics portal (aggregate data, no personal records).\par\vspace{0.4em}
  {\small "The first true AI safety incident." \dimtext{Logan Graham, Anthropic}}
\end{frame}

\begin{frame}{Not just OpenAI}
  \begin{columns}[T,onlytextwidth]
    \column{0.55\textwidth}
      \footnotesize July 30: after the OpenAI news, Anthropic searched 141,006 of its own evaluation runs. In three, a test environment run with the security firm Irregular had real internet access by mistake.\par\vspace{0.4em}
      \begin{itemize}\footnotesize\setlength\itemsep{0.25em}
        \item \textbf{Claude Opus 4.7} saw it was on real production systems and kept going: credentials, a production database
        \item \textbf{Claude Mythos 5} decided it was still in a simulation and published a malicious package to PyPI, which outside systems ran
        \item \textbf{An internal research model} concluded the target was real and stopped
      \end{itemize}
    \column{0.4\textwidth}
      \lbl{Since then}\par\vspace{0.3em}\footnotesize
      Meta (Aug 5) and Google (Sept 18) disclosed test breaches of their own. US lawmakers introduced an AI Kill Switch Act; California's attorney general subpoenaed OpenAI.\par\vspace{0.6em}
      \begin{callout}[warn]{The pattern}
        \footnotesize A test that was meant to be sealed wasn't, and the model treated the real world as part of the task.
      \end{callout}
  \end{columns}
\end{frame}

\begin{frame}{What it means for your setup}
  \footnotesize\renewcommand{\arraystretch}{1.15}
  \begin{tabularx}{\textwidth}{@{}L{4.6cm}Y@{}}
    \toprule
    \hd{What went wrong at the labs} & \hd{Your version, and the fix}\\
    \midrule
    Safety classifiers switched off & \cmd{--dangerously-skip-permissions} only inside a sandbox (\S\,08)\\
    A sandbox with a way out & \cmd{/sandbox} with a domain allowlist; test the wall (\S\,08)\\
    Impossible tasks, no way to quit & let it say "I can't": a CLAUDE.md rule, a \cmd{/goal} budget (\S\,06, \S\,14)\\
    Credentials left in the open & no secrets the agent can read; read-only tokens (\S\,07)\\
    Nobody read the logs & hooks that log and block; transcripts via \kbd{Ctrl+O} (\S\,07, \S\,13)\\
    1,200 agents, one channel & parallel agents multiply mistakes too (\S\,10)\\
    \bottomrule
  \end{tabularx}\par\vspace{0.6em}
  {\small The model did not need bad intent. It needed a goal, a gap, and nobody watching.}
\end{frame}

% =====================================================================
\setsec{A}
\begin{frame}{Appendix: Claude Code and Codex (1/2)}
  \footnotesize\renewcommand{\arraystretch}{1.05}
  \begin{tabularx}{\textwidth}{@{}L{3.1cm}YY@{}}
    \toprule
    \hd{Concept} & \hd{Claude Code} & \hd{Codex}\\
    \midrule
    Instructions & CLAUDE.md, and AGENTS.md natively & AGENTS.md\\
    Default model & Opus 5.5 & GPT-6.1 Sol\\
    Reasoning depth & \texttt{/effort low} to \texttt{max} & \texttt{model\_reasoning\_effort}\\
    Fast mode & \texttt{/fast}, Opus, up to 2.5x & \texttt{/fast on}, \texttt{service\_tier = "fast"}\\
    Approvals & permission modes, \kbd{Shift+Tab} & \texttt{approval\_policy}\\
    OS sandbox & \texttt{/sandbox} & \texttt{sandbox\_mode}\\
    Full bypass & {\scriptsize\texttt{--dangerously-skip-permissions}} & \texttt{--yolo}\\
    Reusable workflows & skills & custom prompts\\
    Lifecycle hooks & settings hooks and mods & hooks\\
    \bottomrule
  \end{tabularx}
\end{frame}

\begin{frame}{Appendix: Claude Code and Codex (2/2)}
  \footnotesize\renewcommand{\arraystretch}{1.05}
  \begin{tabularx}{\textwidth}{@{}L{3.1cm}YY@{}}
    \toprule
    \hd{Concept} & \hd{Claude Code} & \hd{Codex}\\
    \midrule
    Subagents & Agent tool, \texttt{.claude/agents/} & subagents\\
    Session dashboard & agent view, \texttt{claude agents} & agent command center\\
    Peer teams & agent teams, opt-in & none found in the docs index\\
    Worktrees & \texttt{--worktree}, automatic in agent view & \texttt{--worktree}, experimental\\
    Recurring & \texttt{/goal}, \texttt{/loop}, \texttt{/schedule} & scheduled tasks\\
    Headless & \texttt{claude -p} & \texttt{codex exec}\\
    Cloud & claude.ai/code, routines & Codex cloud\\
    Plugins & marketplaces, mods & plugin CLI, remote marketplaces\\
    Switching over & \texttt{/import codex} & -\\
    \bottomrule
  \end{tabularx}\par\vspace{0.4em}
  {\scriptsize\color{muted}Names drift fast on both sides: when this table and the docs disagree, trust the docs.}
\end{frame}

\setsec{}
\begin{frame}[c]
  \begin{columns}[c,onlytextwidth]
    \column{0.3\textwidth}
      \centering\includegraphics[height=5.2cm]{Roberts-Claude-Shannon.jpg}
    \column{0.64\textwidth}
      {\seriffamily\fontsize{38}{40}\selectfont Questions?}\par\vspace{0.7em}
      {\small\color{muted}Claude is named after \textbf{Claude Shannon} (1916-2001), the father of information theory.}\par\vspace{1.1em}
      {\ttfamily\footnotesize\color{dim}
      handout \ra{} claude-codex-workshop-standalone.html\\[0.2em]
      repo \ra{} \href{https://github.com/HaykTarkhanyan/claude-setup}{github.com/HaykTarkhanyan/claude-setup}\\[0.2em]
      docs \ra{} \href{https://code.claude.com/docs}{code.claude.com/docs}\\[0.2em]
      course \ra{} \href{https://academy.claude.com/products/code}{academy.claude.com}}
  \end{columns}
\end{frame}

\end{document}
